% ============================================================================
% MARCO TEÓRICO
% ============================================================================
% Este capítulo presenta la base teórica y conceptual que sustenta el trabajo.
%
% TIPS PARA CITAS (formato APA, según la norma UTEM):
%
%   Cuando el autor forma parte de la narrativa:
%     Según \textcite{clave}, debido a la evolución...
%     Resultado: "Según Wilkinson (2001), debido a la evolución..."
%
%   Cuando el autor NO forma parte de la narrativa:
%     En un estudio reciente \parencite{clave}, se determinó que...
%     Resultado: "En un estudio reciente (Wilkinson, 2001), se determinó que..."
%
%   Cita textual corta (menos de 40 palabras):
%     Como señala el autor, ``texto citado'' \parencite{clave}.
%
%   Cita textual larga (más de 40 palabras):
%     \begin{quote}
%         Texto citado largo... \parencite{clave}
%     \end{quote}
% ============================================================================

\chapter{Marco Teórico}

% --- Párrafo introductorio del capítulo ---
% [PENDIENTE: Breve introducción que presente el propósito y la organización
% de este capítulo al lector.]


% ============================================================================
% SECCIÓN 1: Fundamentos y Arquitectura de los Sistemas ALPR
% ============================================================================
\section{Fundamentos y Arquitectura de los Sistemas ALPR}

El Reconocimiento Automático de Placas Patentes (ALPR, por sus siglas en inglés) 
constituye una tecnología de visión computacional crítica dentro de las 
infraestructuras de Sistemas de Transporte Inteligente (ITS), permitiendo la 
identificación de vehículos en tiempo real mediante la extracción de caracteres 
alfanuméricos a partir de imágenes o secuencias de video 
\parencite{PDFRealTimeLicense2026, kuyunsaDeepLearningEnhancedAutomatic2025, 
larocaAdvancingMultinationalLicense2025}. 
Su implementación es un pilar fundamental en aplicaciones de seguridad pública, 
gestión automatizada de estacionamientos, control de acceso a áreas restringidas 
y aplicación de leyes de tránsito 
\parencite{buleuDeepLearningBasedSystem2025a, xuDeepLearningAlgorithms2026c, 
zherzdevLPRNetLicensePlate2018}. 
La evolución de estos sistemas responde a la necesidad de superar las ineficiencias 
de los modelos de vigilancia manual, los cuales resultan lentos, costosos y altamente 
propensos a errores humanos, limitando la escalabilidad operativa en entornos urbanos 
densos 
\parencite{meesadAdvancedDeepLearning2025a, wijayaEfficientLicensePlate2024, 
xuDeepLearningAlgorithms2026c}.

La eficacia de un sistema ALPR robusto se mide por su capacidad de operar en entornos 
no restringidos (\textit{unconstrained environments}), donde debe procesar datos bajo 
condiciones ambientales adversas y estocásticas 
\parencite{PDFRealTimeLicense2026, xuDeepLearningAlgorithms2026c, 
vargooraniLicensePlateDetection2024}. 
Entre los principales desafíos técnicos se encuentran las variaciones extremas de 
iluminación ---como el deslumbramiento solar o la penumbra nocturna---, las 
obstrucciones por suciedad o agentes climáticos (lluvia, niebla, nieve) y las 
distorsiones geométricas causadas por el movimiento del vehículo o el ángulo de 
captura de la cámara 
\parencite{buleuDeepLearningBasedSystem2025a, mobeenRealTimeAutomaticLicense2026, 
sonnaraEfficientRealtimeLicense2025, xiongLicensePlateRecognition2026}. 
Por ello, la viabilidad de un prototipo operativo depende del cumplimiento de un 
presupuesto de latencia estricto, habitualmente situado en torno a los 100 ms por 
cuadro, asegurando una toma de decisiones inmediata que es indispensable para el 
flujo vehicular en tiempo real 
\parencite{kuyunsaDeepLearningEnhancedAutomatic2025, 
sonnaraEfficientRealtimeLicense2025}.


\subsection{Definición y Evolución: De la Visión Clásica al \textit{Deep Learning}}

La arquitectura de los sistemas ALPR ha experimentado una transformación paradigmática, 
migrando desde métodos deterministas basados en el procesamiento de imágenes tradicional 
hacia redes neuronales altamente parametrizadas 
\parencite{mobeenRealTimeAutomaticLicense2026, xuDeepLearningAlgorithms2026c}. 
Inicialmente, la localización y el reconocimiento dependían de descriptores manuales 
(\textit{handcrafted features}) y heurísticas de bajo nivel, tales como el detector de 
bordes de Canny, la segmentación por color (espacios HSV/YUV) y operaciones morfológicas 
de dilatación para aislar candidatos a placas 
\parencite{PDFRealTimeLicense2026, mobeenRealTimeAutomaticLicense2026, 
xuDeepLearningAlgorithms2026c}. 
En estas etapas tempranas, técnicas como la Transformada de Radon eran fundamentales 
para la corrección de inclinación (\textit{tilt correction}), mientras que la 
clasificación de caracteres se realizaba mediante algoritmos de coincidencia de 
plantillas (\textit{template matching}) o máquinas de vectores de soporte (SVM) 
\parencite{buleuDeepLearningBasedSystem2025a, sonnaraEfficientRealtimeLicense2025}.

A pesar de su eficiencia computacional inicial, estos métodos clásicos presentan una 
fragilidad crítica ante la variabilidad del entorno real, degradando drásticamente su 
precisión frente al desenfoque por movimiento, oclusiones parciales o cambios bruscos 
en la iluminación ambiental 
\parencite{meesadAdvancedDeepLearning2025a, sonnaraEfficientRealtimeLicense2025, 
vargooraniLicensePlateDetection2024}. 
Esta limitación motivó la adopción del Aprendizaje Profundo (\textit{Deep Learning}), 
específicamente mediante Redes Neuronales Convolucionales (CNN), las cuales permiten 
la extracción automatizada de características jerárquicas abstractas directamente desde 
los datos de píxeles \textit{raw}, eliminando la necesidad de ingeniería de 
características manual 
\parencite{PDFRealTimeLicense2026, swatiEfficientDeepLearning2024, 
xuDeepLearningAlgorithms2026c}.

La evolución dentro del aprendizaje profundo se divide en dos enfoques de detección: 
los detectores de dos etapas (ej. Faster R-CNN) y los de una sola etapa (familia YOLO) 
\parencite{PDFRealTimeLicense2026, swatiEfficientDeepLearning2024}. 
Mientras que los modelos de dos etapas priorizan la precisión mediante una Red de 
Propuesta de Regiones (RPN), los modelos de una sola etapa redefinieron la detección 
como un problema de regresión unificado, logrando un balance superior entre latencia 
y exactitud, lo que ha consolidado a arquitecturas como YOLOv8, YOLOv10 y YOLOv12 
como el estándar para aplicaciones de vigilancia y control de acceso en tiempo real 
\parencite{buleuDeepLearningBasedSystem2025a, mobeenRealTimeAutomaticLicense2026, 
sonnaraEfficientRealtimeLicense2025}.

\subsection{Los Pilares del Sistema: Detección (LPD) y Reconocimiento (LPR)}

Un sistema ALPR de alto desempeño se fundamenta en la integración de dos procesos 
computacionales secuenciales y estrechamente acoplados: la Detección de Placas Patentes 
(LPD, por sus siglas en inglés) y el Reconocimiento de Caracteres (LPR, por sus siglas 
en inglés) 
\parencite{larocaAdvancingMultinationalLicense2025, sonnaraEfficientRealtimeLicense2025}. 
La fase LPD constituye la primera etapa crítica y tiene como objetivo la localización 
espacial de la placa dentro de un fotograma mediante la predicción de cajas delimitadoras 
(\textit{bounding boxes}) o polígonos de cuatro esquinas 
\parencite{buleuDeepLearningBasedSystem2025a, sonnaraEfficientRealtimeLicense2025, 
vargooraniLicensePlateDetection2024}. 
La precisión en esta etapa es determinante, dado que cualquier imprecisión en las 
coordenadas estimadas degrada directamente el rendimiento del reconocimiento de 
caracteres subsecuente 
\parencite{suhartoComparativeEvaluationYOLObased2025, swatiEfficientDeepLearning2024, 
xuDeepLearningAlgorithms2026c}.

La fase LPR consiste en la decodificación de la información visual contenida en la 
región detectada para su transformación en una cadena de texto alfanumérico 
\parencite{kuyunsaDeepLearningEnhancedAutomatic2025, wijayaEfficientLicensePlate2024}. 
Históricamente, esta tarea se abordaba mediante un enfoque basado en segmentación, 
el cual intentaba aislar cada carácter individualmente antes de clasificarlo; sin 
embargo, este paradigma presenta una fragilidad intrínseca ante el ruido visual, 
caracteres adheridos y variaciones de iluminación 
\parencite{sonnaraEfficientRealtimeLicense2025, xiongLicensePlateRecognition2026, 
zherzdevLPRNetLicensePlate2018}.

En contraste, el estado del arte ha convergido hacia el enfoque libre de segmentación 
(\textit{segmentation-free}), el cual trata la lectura de la placa como un problema de 
etiquetado secuencial holístico 
\parencite{kuyunsaDeepLearningEnhancedAutomatic2025, sonnaraEfficientRealtimeLicense2025, 
zherzdevLPRNetLicensePlate2018}. 
Esta metodología emplea redes neuronales recurrentes (RNN) o mecanismos de atención 
(\textit{Transformers}) junto a una función de pérdida de Clasificación Temporal 
Conexionista (CTC), permitiendo que el sistema aprenda a reconocer la secuencia completa 
de la patente sin necesidad de pre-segmentar los caracteres manualmente, lo que incrementa 
significativamente la robustez del prototipo en entornos no restringidos 
\parencite{PDFRealTimeLicense2026, kuyunsaDeepLearningEnhancedAutomatic2025, 
mobeenRealTimeAutomaticLicense2026, zherzdevLPRNetLicensePlate2018}.

\subsection{Flujos de Procesamiento: \textit{Pipelines} Secuenciales vs. Sistemas Integrados (\textit{End-to-End})}

La eficiencia operativa de un sistema ALPR en tiempo real depende críticamente de la 
organización lógica de sus componentes, distinguiéndose en la literatura dos enfoques 
arquitectónicos predominantes: los \textit{pipelines} secuenciales multietapa y los 
sistemas integrados extremo a extremo (\textit{end-to-end}) 
\parencite{xuDeepLearningAlgorithms2026c, sonnaraEfficientRealtimeLicense2025}.

El enfoque secuencial o multietapa es el diseño más extendido y se basa en la ejecución 
de modelos independientes para la localización (LPD) y el reconocimiento (LPR) 
\parencite{larocaAdvancingMultinationalLicense2025, mobeenRealTimeAutomaticLicense2026, 
ammarMultiStageDeepLearningBasedVehicle2023}. 
En este flujo, una primera red neuronal detecta la placa y genera un recorte de la 
imagen (\textit{crop}), el cual es preprocesado y enviado a un segundo motor de OCR 
para la extracción de texto 
\parencite{kuyunsaDeepLearningEnhancedAutomatic2025, mobeenRealTimeAutomaticLicense2026}. 
Si bien este paradigma permite una mayor modularidad y facilidad de entrenamiento 
individual, presenta dos desventajas críticas para dispositivos de borde: la redundancia 
paramétrica, causada por la duplicación de extractores de características 
(\textit{backbones}) en cada etapa, y la propagación de errores, donde cualquier 
imprecisión en la detección inicial degrada irreversiblemente la calidad del 
reconocimiento subsecuente 
\parencite{sonnaraEfficientRealtimeLicense2025, xuDeepLearningAlgorithms2026c}.

Frente a estas limitaciones, han surgido los sistemas integrados extremo a extremo de 
un solo paso (\textit{single-pass}), los cuales unifican la detección y el reconocimiento 
en un único grafo computacional 
\parencite{buleuDeepLearningBasedSystem2025a, sonnaraEfficientRealtimeLicense2025, 
zherzdevLPRNetLicensePlate2018}. 
Estas arquitecturas, como el modelo Light-Edge, emplean un \textit{backbone} compartido 
(frecuentemente basado en ResNet o MobileNet) que extrae un mapa de características 
único para ambas tareas, eliminando transferencias de datos ineficientes entre CPU y GPU 
\parencite{sonnaraEfficientRealtimeLicense2025}. 
Al optimizar una función de pérdida conjunta, estos modelos reducen la latencia de 
inferencia de forma drástica ---logrando tasas de procesamiento superiores a los 14 FPS 
en hardware como NVIDIA Jetson Nano--- y minimizan la huella de memoria, consolidándose 
como el estándar de mayor proyección científica para el despliegue de soluciones 
inteligentes en el borde periférico 
\parencite{sonnaraEfficientRealtimeLicense2025, zherzdevLPRNetLicensePlate2018}.


% ============================================================================
% SECCIÓN 2: Análisis del Estado del Arte en Detección (LPD)
% ============================================================================
\section{Análisis del Estado del Arte en Detección (LPD)}

La arquitectura de un sistema ALPR comienza con la capacidad de aislar la región de 
interés (ROI) que contiene la placa patente dentro de una escena compleja. En la 
literatura científica, los algoritmos de detección de objetos se dividen 
taxonómicamente en dos categorías: detectores de dos etapas (\textit{two-stage}) y 
detectores de una sola etapa (\textit{one-stage}) 
\parencite{PDFRealTimeLicense2026, xuDeepLearningAlgorithms2026c}. 
Mientras que los modelos de dos etapas (como Faster R-CNN) utilizan una Red de 
Propuesta de Regiones (RPN) para preseleccionar candidatos, los modelos de una sola 
etapa realizan la clasificación y la regresión de coordenadas en un único paso 
computacional. Esta distinción es crítica para el control de acceso vehicular, ya que 
la naturaleza estocástica del tráfico exige algoritmos que prioricen la eficiencia 
temporal sin sacrificar la robustez espacial 
\parencite{oluwaseyiComparativeAnalysisYOLOv52025, swatiEfficientDeepLearning2024, 
xuDeepLearningAlgorithms2026c}.

\subsection{La Evolución de la Familia YOLO (\textit{You Only Look Once})}

Los modelos de la serie YOLO se han consolidado como el estándar de la industria 
debido a su capacidad para reformular la detección como un problema de regresión 
unificado \parencite{athilakshmiAutomaticNumberPlate2023, mobeenRealTimeAutomaticLicense2026}. 
Su evolución técnica refleja una búsqueda constante por optimizar el flujo de gradientes 
y reducir la latencia de post-procesamiento.

\textbf{YOLOv5}: Representa la madurez de las arquitecturas basadas en cajas de anclaje 
(\textit{anchor-based}). Utiliza descriptores geométricos predefinidos que el modelo emplea 
como referencia para ajustar las coordenadas del objeto 
\parencite{athilakshmiAutomaticNumberPlate2023, oluwaseyiComparativeAnalysisYOLOv52025}. 
Aunque su velocidad es sobresaliente, alcanzando los 92 FPS en entornos controlados, su 
rigidez ante cambios de escala limita su eficacia cuando la cámara de portería captura 
vehículos a diferentes distancias o ángulos oblicuos 
\parencite{oluwaseyiComparativeAnalysisYOLOv52025, swatiEfficientDeepLearning2024}. 
No obstante, su estabilidad la posiciona como un \textit{baseline} confiable para este 
proyecto en caso de requerir una implementación de baja complejidad.

\textbf{YOLOv7}: Introdujo innovaciones en el entrenamiento mediante estrategias de 
``bolsa de obsequios'' (\textit{bag-of-freebies}), optimizando el costo de inferencia sin 
aumentar la carga paramétrica 
(Wang et al., 2023 citado en \cite{mobeenRealTimeAutomaticLicense2026}; \cite{wijayaEfficientLicensePlate2024}). 
Su arquitectura destaca por el uso de redes de agregación de capas eficientes, logrando 
un balance que la hace viable para dispositivos con recursos limitados como la Raspberry Pi, 
permitiendo detecciones precisas en sistemas de portones inteligentes con una latencia controlada 
\parencite{wijayaEfficientLicensePlate2024}.

\textbf{YOLOv8}: Marcó un hito al adoptar un cabezal de detección libre de anclas 
(\textit{anchor-free}), lo que permite predecir directamente el centro de los objetos en 
lugar de ajustar cajas predefinidas, agilizando notablemente la convergencia 
\parencite{mobeenRealTimeAutomaticLicense2026, YOLOv10VsYOLOv82025}. 
Técnicamente, su éxito reside en el módulo C2f (Cross-Stage Partial bottleneck con dos 
convoluciones), el cual combina características de alto nivel con flujos de gradientes 
optimizados, permitiendo al modelo aprender representaciones ricas con una menor 
densidad paramétrica 
\parencite{mobeenRealTimeAutomaticLicense2026, YOLOv10VsYOLOv82025}. 
Esta capacidad de representar detalles de grano fino hace de YOLOv8s un candidato ideal 
para el prototipo UTEM, asegurando la detección de placas bajo condiciones ruidosas.

\textbf{YOLOv10}: Revolucionó el despliegue en el borde mediante la implementación de 
la Doble Asignación Consistente durante el entrenamiento. Este mecanismo emplea 
simultáneamente una estrategia ``uno-a-muchos'' (para proporcionar señales de supervisión 
ricas) y una ``uno-a-uno'' que permite una inferencia nativa libre de Supresión de No 
Máximos (NMS-free) 
\parencite{PDFRealTimeLicense2026, YOLOv8VsYOLOv102025, YOLOv10VsYOLOv82025}. 
Al eliminar el cuello de botella del NMS, que tradicionalmente se ejecuta en CPU, 
YOLOv10 logra una latencia predecible y una eficiencia superior en dispositivos embebidos, 
manteniendo precisiones de hasta 99.16\,\% en escenarios de tráfico denso 
\parencite{meesadAdvancedDeepLearning2025a}.

\textbf{YOLOv11 y YOLOv12}: Representan la integración de mecanismos de atención espacial 
interna. YOLOv12, lanzado en febrero de 2025, emplea arquitecturas centradas en la atención 
de largo alcance, lo que le permite ``limpiar'' el mapa de características y enfocarse en la 
placa incluso ante fondos complejos o degradaciones climáticas 
\parencite{buleuDeepLearningBasedSystem2025a, suhartoComparativeEvaluationYOLObased2025}. 
Resultados experimentales demuestran que YOLOv12 logra un incremento del 2.3\,\% en 
precisión sobre la v11, estableciendo la frontera actual de exactitud para sistemas 
inteligentes de transporte \parencite{buleuDeepLearningBasedSystem2025a}.

\textbf{YOLOv26}: La iteración más reciente (enero 2026) se enfoca en el dominio de la 
inferencia por CPU, logrando ser un 43\,\% más rápida que generaciones previas. Esto se 
consigue mediante la eliminación de la Pérdida Focal de Distribución (DFL), simplificando 
la representación probabilística de las cajas delimitadoras para aliviar la carga de los 
procesadores de bajo costo 
\parencite{YOLOv8VsYOLOv102025, YOLOv10VsYOLOv82025}. 
Además, introduce el optimizador MuSGD, inspirado en los paradigmas de estabilidad de los 
grandes modelos de lenguaje, garantizando que el entrenamiento sea resiliente ante la 
escasez de datos reales \parencite{YOLOv8VsYOLOv102025, YOLOv10VsYOLOv82025}. 
Esta arquitectura representa el estándar de ``futuro asegurado'' para el prototipo propuesto 
en esta investigación.

\subsection{Arquitecturas de Detección Alternativas}

A pesar del dominio de la familia YOLO, la literatura documenta alternativas que ofrecen 
ventajas específicas según el entorno de despliegue:

\textbf{Faster R-CNN}: Es el referente de los detectores de dos etapas. Al utilizar un 
extractor de características (\textit{Backbone}) como Inception V2, logra una precisión 
superior en la detección de patrones visuales extremadamente finos 
\parencite{kuyunsaDeepLearningEnhancedAutomatic2025}. 
Sin embargo, su naturaleza secuencial impone una latencia elevada (aprox. 4.5 ms/imagen 
en GPUs de alto rendimiento), lo que la hace inviable para el control de acceso en tiempo 
real sobre hardware de bajo costo como el que se validará en el Campus Ñuñoa 
\parencite{meesadAdvancedDeepLearning2025a, swatiEfficientDeepLearning2024}.

\textbf{EfficientDet-D0}: Desarrollada para maximizar la eficiencia paramétrica mediante 
redes de pirámide de características bidireccionales (BiFPN), permitiendo una fusión de 
características de diferentes escalas de forma equilibrada 
\parencite{swatiEfficientDeepLearning2024}. 
Aunque su precisión en localización es competitiva, su rendimiento suele ser inferior a 
las versiones recientes de YOLO cuando se enfrenta a distorsiones de perspectiva severas 
en el flujo vehicular continuo \parencite{swatiEfficientDeepLearning2024}.

\textbf{SSD (Single Shot Detector)}: Fue pionera en la detección de un solo paso, 
prediciendo cajas desde múltiples mapas de características. No obstante, estudios 
recientes indican que presenta fluctuaciones de precisión ante objetos pequeños (como 
placas distantes) y una menor capacidad de generalización comparada con los modelos 
\textit{anchor-free} modernos, situándose frecuentemente al final de los \textit{benchmarks} 
de robustez ambiental 
\parencite{meesadAdvancedDeepLearning2025a, swatiEfficientDeepLearning2024}.


% ============================================================================
% SECCIÓN 3: Motores de Extracción de Texto (LPR) y OCR
% ============================================================================
\section{Motores de Extracción de Texto (LPR) y OCR}

Una vez aislada la región de la placa patente, el sistema debe transcribir los patrones 
visuales en cadenas alfanuméricas estructuradas. La eficiencia de este proceso es 
determinante para el rendimiento global del sistema de control de acceso.

\subsection{Metodologías de Reconocimiento}

Históricamente, el reconocimiento de caracteres vehiculares se abordó mediante el enfoque 
basado en segmentación. Este paradigma tradicional fragmenta la placa detectada en 
sub-imágenes que contienen letras y números individuales para ser procesados por 
clasificadores independientes \parencite{xuDeepLearningAlgorithms2026c}. 
Arquitectónicamente, depende de algoritmos de binarización y componentes conectados para 
aislar los glifos, lo que lo hace extremadamente vulnerable a fallos ante caracteres 
adheridos, ruido visual, desenfoque por movimiento o variaciones en la alineación del vehículo 
\parencite{sonnaraEfficientRealtimeLicense2025, xuDeepLearningAlgorithms2026c, zherzdevLPRNetLicensePlate2018}. 
Cualquier error en la segmentación inicial resulta en una falla irreversible en la etapa de 
clasificación \parencite{xiongLicensePlateRecognition2026}.

En contraposición, las arquitecturas modernas adoptan un enfoque libre de segmentación 
(\textit{segmentation-free}), tratando la lectura de la placa como un problema de etiquetado 
secuencial holístico. Este método procesa la imagen completa como una secuencia continua de 
características espaciales, evitando la frágil etapa de binarización manual y recorte 
\parencite{kuyunsaDeepLearningEnhancedAutomatic2025, xuDeepLearningAlgorithms2026c}. 
Al permitir que la red aprenda representaciones directamente de los píxeles sin procesar, 
se incrementa significativamente la robustez en entornos no restringidos 
\parencite{PDFRealTimeLicense2026, sonnaraEfficientRealtimeLicense2025}.

\textbf{Arquitecturas CRNN y el Rol del CTC}: 
El estándar de facto para este enfoque holístico son las Redes Neuronales Recurrentes 
Convolucionales (CRNN). Estas combinan las fortalezas de las redes convolucionales (CNN) 
para la extracción de características con las redes recurrentes (RNN) para el modelado de secuencias 
\parencite{vargooraniLicensePlateDetection2024, zherzdevLPRNetLicensePlate2018}. 
Internamente, una CNN (como ResNet o MobileNet) actúa como \textit{backbone} para generar un 
mapa de características. Luego, capas bidireccionales de Memoria a Corto y Largo Plazo (Bi-LSTM) 
escanean este mapa para capturar dependencias contextuales entre caracteres, entendiendo el 
orden y la estructura sintáctica de la patente 
\parencite{kuyunsaDeepLearningEnhancedAutomatic2025, mobeenRealTimeAutomaticLicense2026, vargooraniLicensePlateDetection2024}.

Para viabilizar el entrenamiento de extremo a extremo sin etiquetas segmentadas a nivel 
de carácter, la capa de Clasificación Temporal Conexionista (CTC) resulta indispensable 
\parencite{sonnaraEfficientRealtimeLicense2025, zherzdevLPRNetLicensePlate2018}. 
El CTC resuelve el problema de alineación prediciendo una distribución de probabilidad para 
cada intervalo espacial, utilizando un carácter ``blanco'' (\textit{blank}) para manejar espacios 
o separar caracteres repetidos \parencite{mobeenRealTimeAutomaticLicense2026, sonnaraEfficientRealtimeLicense2025, zherzdevLPRNetLicensePlate2018}.

\textbf{Transformers y Mecanismos de Atención}: 
La literatura más reciente documenta la evolución hacia mecanismos de atención, sustituyendo 
las capas LSTM por bloques tipo Transformer. Modelos como el Spatial Visual Transformer (SVTR) 
permiten capturar dependencias espaciales de largo alcance, resultando en mapas de 
características más limpios y resilientes ante ruidos o ángulos inclinados 
\parencite{buleuDeepLearningBasedSystem2025a, xuDeepLearningAlgorithms2026c}. 
Incluso, se aprovecha la atención global para mejorar el reconocimiento en capturas 
borrosas de baja resolución (Azadbakht et al., 2022 citado en \cite{xiongLicensePlateRecognition2026}).

\subsection{Benchmarking de Motores y Viabilidad en el Borde}

La elección del motor OCR impacta directamente en la viabilidad técnica del despliegue en 
el Campus Ñuñoa, donde las restricciones de \textit{hardware} exigen un estricto equilibrio 
entre latencia y precisión.

\textbf{Tesseract OCR}: 
Mantenido históricamente por Google, Tesseract utiliza redes LSTM, pero su diseño está 
fuertemente optimizado para documentos impresos estáticos 
\parencite{athilakshmiAutomaticNumberPlate2023, meesadAdvancedDeepLearning2025a, swatiEfficientDeepLearning2024}. 
Su naturaleza secuencial de ejecución en CPU restringe su rendimiento a aproximadamente 
25 lecturas por minuto, haciéndolo ineficiente para flujos de video continuo 
\parencite{xuDeepLearningAlgorithms2026c}. 
Adicionalmente, experimenta una severa degradación de precisión ante distorsiones de 
perspectiva o rotaciones vehiculares \parencite{swatiEfficientDeepLearning2024, xuDeepLearningAlgorithms2026c}.

\textbf{EasyOCR}: 
Esta biblioteca basada en PyTorch emplea un detector de texto CRAFT acoplado a una red 
CRNN con decodificador CTC \parencite{kuyunsaDeepLearningEnhancedAutomatic2025, mobeenRealTimeAutomaticLicense2026, xuDeepLearningAlgorithms2026c}. 
Si bien exhibe una alta robustez ante tipografías complejas y placas deterioradas, su 
carga computacional en entornos carentes de aceleración gráfica es crítica. Su ejecución en CPU 
apenas alcanza 8 lecturas por minuto, mostrando niveles de confianza altamente oscilatorios 
(promedio de 0.41) en secuencias de video real 
\parencite{mobeenRealTimeAutomaticLicense2026, xuDeepLearningAlgorithms2026c}.

\textbf{PaddleOCR}: 
Posicionado como un pipeline de grado industrial, integra el algoritmo de binarización 
DBNet con el modelo de reconocimiento SVTR basado en Transformers 
\parencite{buleuDeepLearningBasedSystem2025a, xuDeepLearningAlgorithms2026c}. 
Es el motor más rápido y preciso documentado, logrando procesar hasta 120 lecturas por 
minuto (en GPU) y alcanzando precisiones del 99.6\,\% al integrarse con versiones recientes 
de YOLO \parencite{buleuDeepLearningBasedSystem2025a, xuDeepLearningAlgorithms2026c}. 
Su manejo eficiente de texto ruidoso y placas oblicuas lo establece como la alternativa 
de código abierto predominante para aplicaciones ALPR \parencite{xuDeepLearningAlgorithms2026c}.

\textbf{Impacto en Hardware de Borde}: 
Para el prototipo UTEM, la NVIDIA Jetson Nano emerge como la plataforma ideal, aprovechando 
sus 128 núcleos CUDA para optimizar arquitecturas como PaddleOCR mediante TensorRT. 
La cuantización del modelo a enteros de 8 bits (INT8) garantiza tasas superiores a los 
14 FPS, satisfaciendo la cuota de latencia de 100\,ms exigida para sistemas de barrera 
automatizada \parencite{sonnaraEfficientRealtimeLicense2025, swatiEfficientDeepLearning2024, xuDeepLearningAlgorithms2026c}. 
En contraste, la utilización de una Raspberry Pi convencional requeriría forzosamente la 
integración de un acelerador NPU dedicado (como el Hailo-8L) para evitar que el cuello 
de botella de inferencia por CPU impida una operación fluida 
\parencite{sonnaraEfficientRealtimeLicense2025, xuDeepLearningAlgorithms2026c}.


% ============================================================================
% SECCIÓN 4: Hardware de Borde y Optimización (Edge AI)
% ============================================================================
\section{Hardware de Borde y Optimización (\textit{Edge AI})}

La transición de modelos desde servidores de alto rendimiento hacia el borde de la 
red (\textit{Edge AI}) es un requerimiento fundamental para sistemas ALPR operativos, 
donde la dependencia de la nube introduce latencias inaceptables y riesgos de privacidad. 
Sin embargo, el despliegue local impone severas restricciones de memoria, consumo 
energético y capacidad de cómputo.

\subsection{Descentralización y Sincronización Local (\textit{Edge Caching})}

Uno de los desafíos principales en arquitecturas IoT orientadas al control de acceso es mitigar la dependencia de la red troncal para la validación de los datos. Si bien los dispositivos de borde realizan la inferencia de visión artificial localmente, la decisión final (autorizar una patente) requiere contrastar el resultado contra una base de datos. Realizar esta validación de forma síncrona contra la nube reintroduce latencias críticas de red y expone el sistema a fallos por desconexión. La literatura enfatiza que la ventaja del despliegue en el borde es precisamente ``cortar el vínculo entre internet y la máquina host, evitando así la latencia de la red y creando un tiempo de respuesta más rápido'' [traducción propia] \parencite{swatiEfficientDeepLearning2024}.

Para resolver este desafío de diseño, el estado del arte en redes IoT incorpora el patrón arquitectónico de descentralización de datos o \textit{Edge Caching}. Esta estrategia consiste en almacenar la información de las patentes habilitadas en una base de datos local dentro del dispositivo \parencite{xuDeepLearningAlgorithms2026c}. El nodo perimetral sincroniza periódicamente la lista maestra en segundo plano, logrando que la consulta sea local e inmediata. Esto fomenta un entorno de procesamiento distribuido que reduce drásticamente el consumo de ancho de banda y mejora la escalabilidad \parencite{buleuDeepLearningBasedSystem2025a, kuyunsaDeepLearningEnhancedAutomatic2025}. Así, el sistema garantiza resiliencia operativa (\textit{offline-first}), asegurando que el control de acceso actúe en el orden de los milisegundos sin depender de servidores centrales.

\subsection{Ecosistemas de Cómputo en el Borde: Arquitecturas y Rendimiento}

La literatura documenta dos grandes familias de ecosistemas de cómputo para el despliegue 
de redes neuronales profundas en el borde: aquellas con aceleración nativa por GPU y 
aquellas basadas primordialmente en procesamiento de CPU.

\textbf{Aceleración Nativa por GPU (Arquitectura NVIDIA Jetson)}: 
La familia NVIDIA Jetson (incluyendo Nano, Xavier AGX y Orin) está diseñada específicamente 
para satisfacer las demandas de cómputo paralelo masivo del \textit{Deep Learning} 
\parencite{ammarMultiStageDeepLearningBasedVehicle2023, sonnaraEfficientRealtimeLicense2025}. 
A diferencia de las CPUs convencionales, la Jetson Nano integra una GPU de arquitectura Maxwell 
con 128 núcleos CUDA, lo que permite la ejecución concurrente de operaciones tensoriales. 
Esta paralelización optimiza el rendimiento en modelos de detección de una sola etapa como 
YOLOv8 y YOLOv10 \parencite{meesadAdvancedDeepLearning2025a, sonnaraEfficientRealtimeLicense2025}. 
El ecosistema se apoya en la suite NVIDIA JetPack, la cual incluye el motor de inferencia 
TensorRT, capaz de optimizar los grafos computacionales para maximizar la tasa de procesamiento 
(FPS) y cumplir con el presupuesto de latencia de $\approx$100\,ms, métrica exigida para 
el control de acceso vehicular en tiempo real \parencite{ammarMultiStageDeepLearningBasedVehicle2023, sonnaraEfficientRealtimeLicense2025}.

\textbf{Procesamiento Basado en CPU (Arquitectura Raspberry Pi)}: 
Por otro lado, plataformas como la Raspberry Pi 5 utilizan procesadores de propósito 
general (e.g., Broadcom BCM2712 Cortex-A76) que, si bien son eficientes en tareas 
lógicas secuenciales, carecen de unidades de procesamiento gráfico compatibles con 
marcos de IA de alto nivel como CUDA \parencite{oluwaseyiComparativeAnalysisYOLOv52025}. 
Ejecutar modelos de detección complejos directamente en su CPU resulta en una latencia 
prohibitiva para aplicaciones vehiculares, logrando frecuentemente tasas de inferencia 
inferiores a 5 FPS \parencite{oluwaseyiComparativeAnalysisYOLOv52025, xuDeepLearningAlgorithms2026c}. 
Esto podría traducirse en la pérdida de detección de vehículos con movimiento rápido en 
entornos de portería.

Para viabilizar el uso de placas como la Raspberry Pi en tareas de inferencia pesada, 
resulta imperativa la incorporación de Unidades de Procesamiento Neuronal (NPU) externas, 
como el módulo AI HAT+ basado en chips Hailo-8L \parencite{oluwaseyiComparativeAnalysisYOLOv52025}. 
Arquitectónicamente, estas NPUs están optimizadas para los flujos de datos de redes neuronales, 
ejecutando cálculos matriciales con una eficiencia energética (TOPS por Watt) muy superior 
a la de una CPU convencional \parencite{oluwaseyiComparativeAnalysisYOLOv52025, xuDeepLearningAlgorithms2026c}.

\subsection{Técnicas de Compilación y Cuantización}

\textbf{Fundamentos de la Cuantización}: 
Para adaptar modelos masivos a las restricciones de memoria del hardware de borde, se 
emplea el proceso de cuantización. Esta técnica consiste en reducir la precisión numérica 
de los pesos y las funciones de activación de la red neuronal, trasladándolos desde 
representaciones de punto flotante de 32 bits (FP32) a formatos más compactos como 16 bits 
(FP16) o números enteros de 8 bits (INT8) 
\parencite{sonnaraEfficientRealtimeLicense2025, swatiEfficientDeepLearning2024}. 
Matemáticamente, esto implica mapear los valores continuos del modelo original a un 
conjunto discreto de niveles, utilizando frecuentemente la calibración por divergencia 
de Kullback-Leibler (KL) para minimizar la distorsión estadística de los datos 
\parencite{ammarMultiStageDeepLearningBasedVehicle2023}.

\textbf{Beneficios en Memoria y Aceleración de Inferencia}: 
La transición de FP32 a INT8 reduce la huella del modelo en disco y RAM en un factor 
de hasta 4x. En contextos operativos, modelos ALPR con un peso inicial de 37 MB logran 
comprimirse a 11 MB al ser cuantizados a 16 bits, un factor crítico para microordenadores 
con recursos limitados (e.g., los 4 GB de RAM de la Jetson Nano) 
\parencite{sonnaraEfficientRealtimeLicense2025, swatiEfficientDeepLearning2024}. 
A nivel de ejecución, el menor tamaño de los datos disminuye drásticamente el ancho de 
banda requerido para las transferencias de memoria. Además, el hardware moderno incorpora 
instrucciones SIMD que permiten procesar múltiples operaciones INT8 en un único ciclo de 
reloj, posibilitando aceleraciones de inferencia de hasta 4.6x 
\parencite{ammarMultiStageDeepLearningBasedVehicle2023, sonnaraEfficientRealtimeLicense2025}.

\textbf{Rol de Motores de Optimización}: 
La eficiencia en el despliegue requiere de compiladores especializados. 
TensorRT, por ejemplo, ejecuta una "fusión de capas" (\textit{Layer \& Tensor Fusion}), 
combinando operaciones de convolución, sesgo y activación en un único \textit{kernel} de 
cómputo, lo que mitiga las ineficientes transferencias de datos en memoria 
\parencite{ammarMultiStageDeepLearningBasedVehicle2023, sonnaraEfficientRealtimeLicense2025}. 
Por su parte, TensorFlow Lite ofrece formatos de serialización (\textit{FlatBuffers}) e 
intérpretes optimizados para arquitecturas ARM, facilitando el despliegue en microprocesadores 
móviles \parencite{swatiEfficientDeepLearning2024}.

\textbf{Impacto en la Precisión (mAP)}: 
Pese a la compresión numérica, la literatura empírica demuestra que la pérdida de 
precisión media (mAP) es marginal cuando se emplean técnicas de calibración rigurosas. 
En evaluaciones del modelo Light-Edge optimizado mediante Torch-TensorRT a INT8, se 
reportó un incremento de velocidad de 3.1 FPS a 14.2 FPS, incurriendo en una penalización 
de apenas 0.4 puntos porcentuales en mAP respecto al modelo FP32 base 
\parencite{sonnaraEfficientRealtimeLicense2025}. 
Esto evidencia que la cuantización no es una concesión comprometedora, sino la técnica 
habilitadora que hace factible el funcionamiento de arquitecturas de \textit{Deep Learning} 
en escenarios comerciales de borde 
\parencite{sonnaraEfficientRealtimeLicense2025, swatiEfficientDeepLearning2024}.


% ============================================================================
% SECCIÓN 5: Robustez en Entornos No Restringidos
% ============================================================================
\section{Robustez en Entornos No Restringidos}

El despliegue de un sistema ALPR en exteriores, como las porterías vehiculares del 
Campus Ñuñoa, somete a los algoritmos a una drástica variabilidad ambiental. 
Garantizar la resiliencia del sistema ante escenarios estocásticos (fluctuaciones 
lumínicas, sombras dinámicas y precipitaciones) requiere la integración de módulos 
especializados a lo largo de todo el flujo de procesamiento.

\subsection{Técnicas de Preprocesamiento Crítico}

Aunque las arquitecturas de aprendizaje profundo poseen una capacidad intrínseca para 
la extracción autónoma de características, el preprocesamiento de la imagen sigue 
siendo una etapa indispensable. Depurar las variaciones estocásticas antes de la 
inferencia reduce la complejidad requerida en el \textit{backbone} de la red, 
aliviando significativamente la carga computacional en los dispositivos de borde 
\parencite{kuyunsaDeepLearningEnhancedAutomatic2025, xiongLicensePlateRecognition2026}.

Entre las técnicas fundamentales, la Normalización Min-Max ($x_{norm} = \frac{x - x_{min}}{x_{max} - x_{min}}$) 
es vital para estabilizar el entrenamiento. Al escalar las intensidades de los píxeles, 
evita que valores extremos dominen el cálculo de los pesos, acelerando la convergencia 
del descenso de gradiente \parencite{kuyunsaDeepLearningEnhancedAutomatic2025}. 
Por otra parte, para el manejo de sombras o deslumbramiento severo, la Ecualización 
de Histograma Adaptativa Limitada por Contraste (CLAHE) demuestra superioridad sobre la 
ecualización global. Al operar mediante regiones localizadas (\textit{tiles}), CLAHE 
optimiza el contraste sin amplificar el ruido en zonas homogéneas, permitiendo que 
detectores como YOLOv11 aíslen la placa con mayor fiabilidad 
\parencite{buleuDeepLearningBasedSystem2025a, suhartoComparativeEvaluationYOLObased2025}.

Para la atenuación de ruido visual sin comprometer la morfología alfanumérica, la 
literatura prioriza los Filtros Bilaterales sobre el filtrado Gaussiano tradicional. 
Esta técnica logra suavizar el granulado digital mientras preserva los bordes de alta 
frecuencia, evitando confusiones semánticas en la etapa OCR (e.g., confundir una ``B'' 
con un ``8'') \parencite{kuyunsaDeepLearningEnhancedAutomatic2025, xiongLicensePlateRecognition2026}. 
Finalmente, transformaciones como la conversión a escala de grises seguida de una 
binarización de Otsu maximizan la dicotomía entre fondo y texto, reduciendo drásticamente 
la dimensionalidad de los datos de entrada 
\parencite{kuyunsaDeepLearningEnhancedAutomatic2025, suhartoComparativeEvaluationYOLObased2025}.

\subsection{Módulos de Mejora Dinámica: Restauración y Adaptación}

Las perturbaciones climáticas severas exigen enfoques más sofisticados que el preprocesamiento 
estadístico clásico. Recientemente, se han propuesto arquitecturas con módulos de mejora 
dinámica diseñados explícitamente para restaurar características degradadas por la niebla 
o la lluvia.

Destaca el módulo CPA-Enhancer (\textit{Chain-of-Thought Prompted Adaptive Enhancer}), un bloque 
que se acopla previo al \textit{backbone} de detección \parencite{xiongLicensePlateRecognition2026}. 
Utilizando convoluciones de atención de campo receptivo (RFAConv), extrae información 
contextual para generar representaciones optimizadas a múltiples resoluciones. La fusión de 
estas señales reemplaza la imagen degradada original, logrando incrementos sustanciales de 
precisión en escenarios climáticos adversos \parencite{xiongLicensePlateRecognition2026}.

En paralelo, para abordar la baja resolución provocada por la lejanía del vehículo, 
módulos como CARAFE (\textit{Content-Aware ReAssembly of FEatures}) reemplazan algoritmos 
clásicos de sobremuestreo (como la interpolación bilineal). CARAFE reconstruye los trazos 
de los caracteres adaptándose dinámicamente al contenido visual 
\parencite{xiongLicensePlateRecognition2026}. De forma similar, el \textit{framework} SR2 
(Super-Resolución y Reconocimiento) aborda este problema entrenando ambas tareas de manera 
paralela, utilizando la super-resolución como un regulador inductivo para refinar los 
mapas de características que ingresan al motor de clasificación 
\parencite{schirrmacherBenchmarkingProbabilisticDeep2023}.

Finalmente, la integración de mecanismos de atención espacial y de canal (como SEAM o CBAM) 
permite a la red asignar ponderaciones dinámicas a las regiones de interés, ignorando 
la lluvia de fondo para enfocarse estrictamente en los glifos de la matrícula 
\parencite{xiongLicensePlateRecognition2026, xuDeepLearningAlgorithms2026c}.

\subsection{Datos Sintéticos y Aumento (\textit{Data Augmentation})}

Uno de los principales impedimentos en la investigación de ALPR es la escasez de conjuntos de 
datos reales anotados bajo condiciones extremas. Para solventar este déficit y evitar el 
sobreajuste (\textit{overfitting}), la investigación moderna recurre a la expansión controlada 
del volumen de entrenamiento \parencite{larocaAdvancingMultinationalLicense2025, meesadAdvancedDeepLearning2025a}.

Las técnicas de aumento de datos (\textit{Data Augmentation}) en tiempo real (\textit{on-the-fly}) 
constituyen la base de este proceso. Alteraciones estocásticas como rotaciones ($\pm15^\circ$), 
recortes aleatorios, ajustes en el espacio de color HSV y transformaciones de mosaico, 
simulan oclusiones parciales y diversas perspectivas angulares 
\parencite{buleuDeepLearningBasedSystem2025a, vargooraniLicensePlateDetection2024}. 
Complementariamente, la Permutación de Caracteres se utiliza para generar variabilidad 
léxica, intercambiando parches de caracteres dentro de una misma matrícula sin alterar 
la iluminación o el ruido original de la imagen \parencite{larocaAdvancingMultinationalLicense2025}.

Cuando las transformaciones geométricas son insuficientes, las Redes Generativas 
Antagónicas (GANs) asumen el rol de síntesis avanzada. Arquitecturas como \textit{pix2pix} o 
LPSRGAN son capaces de realizar traducción de dominios; por ejemplo, convirtiendo imágenes 
diurnas limpias en capturas nocturnas con ruido o introduciendo degradaciones hiperrealistas 
como lodo salpicado \parencite{larocaAdvancingMultinationalLicense2025, xuDeepLearningAlgorithms2026c}. 
La inclusión de esta data sintética de alta fidelidad permite entrenar sistemas altamente 
robustos incluso empleando menos del 10\,\% de imágenes reales, un enfoque metodológico 
crucial para generalizar el desempeño del prototipo 
\parencite{larocaAdvancingMultinationalLicense2025}.


% ============================================================================
% SECCIÓN 6: Consistencia Temporal y Seguimiento (Tracking)
% ============================================================================
\section{Consistencia Temporal y Seguimiento (\textit{Tracking})}

La detección estática en un único fotograma es insuficiente para un entorno operativo 
como el control de acceso vehicular. El sistema debe garantizar la continuidad de la 
identidad del vehículo a través del tiempo para evitar registrar el mismo evento de 
acceso de forma redundante \parencite{ammarMultiStageDeepLearningBasedVehicle2023, mobeenRealTimeAutomaticLicense2026}. 
La integración de la dimensión temporal transforma detecciones aisladas en un sistema 
robusto y cohesivo.

\subsection{Algoritmos de Seguimiento (\textit{Tracking})}

Los algoritmos de \textit{tracking} vinculan las detecciones realizadas en el fotograma 
$k$ con las identidades confirmadas en fotogramas previos, imponiendo una lógica 
espacio-temporal que mitiga los fallos transitorios de los modelos de visión artificial 
\parencite{mobeenRealTimeAutomaticLicense2026}.

En entornos con limitaciones de recursos, el algoritmo SORT (\textit{Simple Online and Realtime Tracking}) 
se perfila como un estándar de alta eficiencia computacional. Arquitectónicamente, SORT emplea un 
Filtro de Kalman para realizar predicciones lineales del estado cinemático del vehículo 
(posición, escala y velocidad). Posteriormente, utiliza el Algoritmo Húngaro para la 
asociación de datos, minimizando la métrica de Intersección sobre la Unión (IoU) entre 
las cajas delimitadoras predichas y las recientemente detectadas 
\parencite{mobeenRealTimeAutomaticLicense2026}.

Ante oclusiones severas o intersecciones visuales entre múltiples vehículos, algoritmos 
más avanzados como DeepSORT incorporan una métrica de asociación profunda. Mediante la 
extracción de descriptores visuales (\textit{embeddings}) a través de redes neuronales, 
DeepSORT asegura que el seguimiento no dependa exclusivamente de la geometría del 
movimiento, sino también de la apariencia física del vehículo, minimizando el 
intercambio erróneo de identidades (\textit{ID switches}) 
\parencite{ammarMultiStageDeepLearningBasedVehicle2023, mobeenRealTimeAutomaticLicense2026}.

Mediante estos algoritmos, se asigna un identificador único (\texttt{track\_id}) a 
cada vehículo. Esto permite que el nodo de borde acumule información internamente y 
envíe una única notificación consolidada al panel administrativo, reduciendo 
drásticamente la saturación de red y asegurando la trazabilidad del evento 
\parencite{ammarMultiStageDeepLearningBasedVehicle2023}.

\subsection{Votación Temporal e Interpolación}

El procesamiento de secuencias de video permite aplicar mecanismos de Votación Temporal 
(\textit{Majority Voting}), los cuales aprovechan la redundancia fotogramétrica para 
incrementar la precisión del Reconocimiento de Caracteres (OCR). 

En este paradigma, el sistema almacena las predicciones obtenidas de múltiples 
fotogramas en un búfer temporal vinculado al \texttt{track\_id}. En lugar de confiar 
en la lectura estática ---altamente vulnerable al desenfoque por movimiento o reflejos 
momentáneos---, el sistema efectúa una votación carácter por carácter. El glifo que 
presenta mayor frecuencia de aparición en la secuencia se asume como el valor definitivo 
\parencite{ammarMultiStageDeepLearningBasedVehicle2023}. Variantes más robustas aplican 
una votación ponderada, multiplicando cada voto por el puntaje de confianza (\textit{confidence score}) 
emitido por el motor OCR, otorgando mayor peso analítico a los fotogramas donde el 
vehículo se encuentra más cercano y nítido 
\parencite{ammarMultiStageDeepLearningBasedVehicle2023}. La evidencia empírica documenta que 
esta técnica eleva la precisión final del reconocimiento hasta en un 40\,\% respecto al 
procesamiento de imágenes aisladas \parencite{ammarMultiStageDeepLearningBasedVehicle2023, mobeenRealTimeAutomaticLicense2026}.

Adicionalmente, cuando el detector pierde el rastro de la placa por bloqueos temporales, 
se aplica la interpolación de datos geométricos. Si se registran posiciones válidas en 
los tiempos $t_0$ y $t_1$, una interpolación lineal 1-D permite imputar las coordenadas 
faltantes, reconstruyendo una trayectoria fluida \parencite{mobeenRealTimeAutomaticLicense2026}. 
Esta recuperación sintética de la continuidad visual puede aumentar la densidad de los 
datos de seguimiento en más de un 100\,\%, garantizando que el sistema mantenga el foco 
en la placa bajo condiciones operativas inestables \parencite{mobeenRealTimeAutomaticLicense2026}.


% ============================================================================
% SECCIÓN 7: Contexto Normativo y Operacional Chileno
% ============================================================================
\section{Contexto Normativo y Operacional Chileno}

La viabilidad de un sistema ALPR en el Campus Ñuñoa no solo depende de su arquitectura 
computacional, sino de su estricta alineación con el ecosistema vehicular chileno y 
las regulaciones de privacidad institucionales.

\subsection{Estándares de Patentes en Chile}

El parque automotriz chileno presenta una notable heterogeneidad derivada de la 
evolución de sus normativas de tránsito. Desde el formato clásico instaurado en 1985 
(dos letras y cuatro números), el país evolucionó en 2007 hacia el estándar FE-Schrift 
(cuatro letras y dos números), una tipografía diseñada específicamente para dificultar 
la falsificación y facilitar el reconocimiento óptico \parencite{https://magnet.clComoSeranCuando}. 
Recientemente, se ha aprobado la transición hacia configuraciones de cinco letras y 
un número para expandir las combinaciones disponibles \parencite{https://magnet.clComoSeranCuando}.

Un hito crítico para la visión computacional es la promulgación del Decreto Supremo 53, 
que introduce la ``Patente Verde'' para vehículos de electromovilidad (eléctricos e 
híbridos enchufables). A diferencia del fondo blanco tradicional, este formato emplea 
caracteres negros sobre un fondo verde reflectante 
\parencite{https://magnet.clComoSeranCuando}. Esta variabilidad de fondos y densidades 
alfanuméricas representa un desafío crítico: los sistemas ALPR basados en heurísticas de 
procesamiento clásico sufren drásticas caídas de precisión ante cambios repentinos de 
contraste \parencite{xuDeepLearningAlgorithms2026c}.

Para absorber esta diversidad normativa sin comprometer la eficacia, resulta imperativo 
el uso de modelos de \textit{Deep Learning} como YOLOv8 \parencite{oluwaseyiComparativeAnalysisYOLOv52025, mobeenRealTimeAutomaticLicense2026}. 
Esta arquitectura extrae características jerárquicas que son invariantes al color del 
fondo, asegurando que el control de acceso de la universidad sea inclusivo ante las 
nuevas tecnologías vehiculares \parencite{larocaAdvancingMultinationalLicense2025}. 
Cabe destacar que, dado que las cámaras de portería efectúan un escaneo frontal, 
el sistema excluye inherentemente a las motocicletas, las cuales por normativa chilena 
solo portan placa trasera. Del mismo modo, se excluyen vehículos de carga y emergencia 
cuyos formatos escapan al perfil del funcionario universitario.

\subsection{Requerimientos Funcionales y Trazabilidad}

El diseño lógico del sistema exige el cumplimiento de parámetros operativos rigurosos 
acordes a la filosofía de \textit{Edge Computing}. En primer lugar, para que una 
decisión de acceso sea percibida como inmediata, la literatura impone un presupuesto de 
latencia estricto de $\approx$100\,ms por cuadro procesado en el hardware local 
\parencite{ammarMultiStageDeepLearningBasedVehicle2023, sonnaraEfficientRealtimeLicense2025}. 
Para garantizar esta fluidez durante las ventanas de mayor tráfico en la portería, el 
nodo de borde debe sostener tasas de inferencia objetivo entre 10 y 14 FPS 
\parencite{sonnaraEfficientRealtimeLicense2025}. Mantenerse en este umbral evita que 
un vehículo en movimiento abandone el campo visual antes de que su patente sea 
decodificada con un alto nivel de confianza.

En el ámbito administrativo y normativo, la automatización del acceso debe respaldarse 
mediante la generación de registros digitales históricos (\textit{logs}). Cada evento 
de apertura debe quedar almacenado localmente con una marca de tiempo precisa, la 
cadena alfanumérica extraída y el nivel de confianza arrojado por el modelo 
\parencite{ammarMultiStageDeepLearningBasedVehicle2023}. Esta trazabilidad digital 
sustituye la falta de métricas del modelo manual vigente, permitiendo auditorías 
posteriores y facilitando la planificación estratégica de la capacidad del 
estacionamiento \parencite{kumarAutomaticNumberPlate2026}. Además, en cumplimiento con 
la Ley 19.628 sobre Protección de la Vida Privada \parencite{nacionalBibliotecaCongresoNacional1999}, 
el procesamiento íntegro en el borde (sin envío de secuencias de video a la nube) garantiza la 
confidencialidad de los datos vehiculares de la comunidad universitaria.
